%% ============================================================
\section{Implementation}
\label{sec:implementation}
%% ============================================================

\subsection{Technology Stack Overview}

\echoray{} is a full-stack TypeScript monorepo with two independent
deployable units: the \textbf{client} (Next.js~15 application) and the
\textbf{server} (Express~5 + Socket.io~4 API). All production dependency
versions are pinned in \texttt{package.json}.

\subsubsection{Exact Production Dependency Versions}

\begin{table}[H]
  \centering
  \caption{Pinned production dependencies for \echoray{} (March 2026).}
  \label{tab:deps}
  \renewcommand{\arraystretch}{1.05}
  \begin{tabular}{llp{6cm}}
    \toprule
    \textbf{Package} & \textbf{Version} & \textbf{Purpose} \\
    \midrule
    \multicolumn{3}{l}{\textit{Client (Next.js application)}} \\
    \midrule
    \texttt{next}                    & 15.3.3  & React SSR framework, App Router \\
    \texttt{react}                   & \^{}19.0.0 & UI rendering \\
    \texttt{@monaco-editor/react}    & \^{}4.7.0 & Browser-native VS Code editor \\
    \texttt{@reduxjs/toolkit}        & \^{}2.8.2 & Predictable state management \\
    \texttt{react-redux}             & \^{}9.2.0 & Redux React bindings \\
    \texttt{@webcontainer/api}       & \^{}1.6.1 & Node.js WASM runtime \\
    \texttt{socket.io-client}        & \^{}4.8.1 & WebSocket client \\
    \texttt{axios}                   & \^{}1.9.0 & HTTP client with interceptors \\
    \texttt{gsap}                    & \^{}3.13.0 & 60\,fps animation engine \\
    \texttt{react-markdown}          & \^{}10.1.0 & AI response markdown rendering \\
    \texttt{remark-gfm}              & \^{}4.0.1 & GFM (tables, strikethrough) \\
    \texttt{dompurify}               & \^{}3.2.6 & XSS-safe HTML sanitisation \\
    \texttt{react-hot-toast}         & \^{}2.5.2 & Non-blocking notifications \\
    \texttt{tailwindcss}             & \^{}4.0   & Utility-first CSS \\
    \midrule
    \multicolumn{3}{l}{\textit{Server (Express API + Socket.io)}} \\
    \midrule
    \texttt{express}                 & \^{}5.1.0  & HTTP framework (major upgrade to v5) \\
    \texttt{socket.io}               & \^{}4.8.1  & WebSocket server \\
    \texttt{@google/generative-ai}   & \^{}0.24.1 & Gemini SDK \\
    \texttt{@prisma/client}          & \^{}6.10.1 & ORM with type-safe queries \\
    \texttt{@prisma/extension-accelerate} & \^{}1.3.0 & Edge-cached query layer \\
    \texttt{bcrypt}                  & \^{}5.1.1  & Password hashing \\
    \texttt{jsonwebtoken}            & \^{}9.0.2  & JWT sign/verify \\
    \texttt{ioredis}                 & \^{}5.6.1  & Redis client (token blacklist) \\
    \texttt{zod}                     & \^{}3.25.67 & Runtime input validation \\
    \texttt{helmet}                  & \^{}8.1.0  & HTTP security headers \\
    \texttt{cookie-parser}           & \^{}1.4.7  & Cookie middleware \\
    \texttt{morgan}                  & \^{}1.10.0 & HTTP request logging \\
    \texttt{typescript}              & \^{}5.8.3  & Type-safe JavaScript superset \\
    \texttt{prisma}                  & \^{}6.9.0  & Schema + migration tooling \\
    \bottomrule
  \end{tabular}
\end{table}

\subsubsection{Hardware and Infrastructure Requirements}

\begin{table}[H]
  \centering
  \caption{Minimum, recommended, and production infrastructure requirements.}
  \label{tab:hw}
  \begin{tabular}{lccc}
    \toprule
    \textbf{Component} & \textbf{Minimum} & \textbf{Recommended} & \textbf{Production} \\
    \midrule
    Server CPU  & 0.5 vCPU & 2 vCPU & 4+ vCPU \\
    Server RAM  & 512\,MB  & 2\,GB  & 4\,GB \\
    PostgreSQL  & shared   & 2\,GB RAM & Managed (Supabase/RDS) \\
    Redis       & 25\,MB   & 256\,MB & Upstash/ElastiCache \\
    Client CPU  & Any modern browser & --- & --- \\
    Client WASM & SharedArrayBuffer support required & --- & --- \\
    Network     & 100\,Mbps & 1\,Gbps & CDN + load balancer \\
    \bottomrule
  \end{tabular}
\end{table}

\textbf{WebContainer WASM requirements.} The WebContainer API requires the
following browser headers, which must be configured on the server serving the
HTML:
\begin{center}
\texttt{Cross-Origin-Opener-Policy: same-origin}\\
\texttt{Cross-Origin-Embedder-Policy: require-corp}
\end{center}
These enable \texttt{SharedArrayBuffer}, which WebAssembly threads require.

\subsection{Golden Snippet: The AI-Broadcast Core}

The most critical innovation in \echoray{} is the tight coupling of the
AI service, Socket.io broadcast, and file-tree mount in a single asynchronous
workflow. The \textbf{Golden Snippet}---the minimal code expressing this
innovation---spans three files.

\subsubsection{Server-Side AI Invocation and Broadcast}

\begin{lstlisting}[
  caption={Golden Snippet Part 1: Socket.io event handler with AI routing
    (\texttt{server/src/socket.ts}).},
  label=lst:golden_socket,
  language=Java,
]
socket.on("project-message", async (data: any) => {
  const message = data.message;
  const aiIsPresentInMessage = message.includes("@ai");

  // 1. Always rebroadcast the user's raw message first
  //    (so peers see the query before the AI responds)
  socket.broadcast
        .to(socket.roomId!)
        .emit("project-message", data);

  // 2. If @ai prefix detected, invoke Gemini asynchronously
  if (aiIsPresentInMessage) {
    const prompt = message.replace("@ai", "");          // strip sentinel
    const result = await generateResult(prompt);        // Gemini call

    // 3. Broadcast AI response to ALL room members (including invoker)
    io.to(socket.roomId!).emit("project-message", {
      message: result,
      sender: { _id: "ai", email: "AI" },
    });
  }
});
\end{lstlisting}

\subsubsection{AI Service: Gemini Configuration}

\begin{lstlisting}[
  caption={Golden Snippet Part 2: Gemini model configuration with JSON mode
    and expert system instruction (\texttt{server/src/services/ai.service.ts}).},
  label=lst:golden_ai,
  language=Java,
]
const genAI = new GoogleGenerativeAI(process.env.GEMINI_API_KEY!);

const model = genAI.getGenerativeModel({
  model: "gemini-2.5-flash",
  generationConfig: {
    responseMimeType: "application/json",   // enforce JSON output schema
    temperature: 0.4,                       // deterministic code generation
  },
  systemInstruction: `You are an expert JavaScript/MERN developer.
    Always respond in JSON:
    {
      "text": "<explanation>",
      "fileTree": {
        "<filename>": { "file": { "contents": "<code>" } }
      },
      "buildCommand": { "mainItem": "npm", "commands": ["install"] },
      "startCommand": { "mainItem": "node",  "commands": ["app.js"]  }
    }`,
});

export const generateResult = async (prompt: any): Promise<string> => {
  const result = await model.generateContent(prompt);
  return result.response.text();                        // raw JSON string
};
\end{lstlisting}

\subsubsection{Client-Side: WebContainer Mount from AI Response}

\begin{lstlisting}[
  caption={Golden Snippet Part 3: Parsing AI response and mounting generated
    file tree into the WebContainer runtime
    (\texttt{client/app/project/[id]/page.tsx}).},
  label=lst:golden_wc,
  language=Java,
]
// On receiving an AI message from the socket:
const handleAIMessage = async (rawResult: string) => {
  const parsed = JSON.parse(rawResult);               // { text, fileTree }

  // Display the explanation in the chat
  setMessages(prev => [...prev, {
    message: parsed.text,
    sender: { _id: "ai", email: "AI" }
  }]);

  // If the AI generated code, merge it into the project file tree
  if (parsed.fileTree) {
    const mergedTree = { ...fileTree, ...parsed.fileTree };
    setFileTree(mergedTree);                           // update React state

    // Mount directly into the live WebContainer instance
    const wc = await getWebContainerInstance();
    await wc.mount(parsed.fileTree);                  // zero-copy WASM mount

    // Install dependencies if a buildCommand is present
    if (parsed.buildCommand) {
      const installProc = await wc.spawn(
        parsed.buildCommand.mainItem,
        parsed.buildCommand.commands
      );
      await installProc.exit;
    }
  }
};
\end{lstlisting}

\subsection{Database Schema}

The Prisma schema defines two models with a many-to-many self-join:

\begin{lstlisting}[
  caption={Prisma schema (\texttt{server/prisma/schema.prisma}).},
  label=lst:schema,
  language=Java,
]
model User {
  id        Int       @id @default(autoincrement())
  email     String    @unique
  password  String
  createdAt DateTime  @default(now())
  projects  Project[] @relation("ProjectUsers")
}

model Project {
  id        Int      @id @default(autoincrement())
  name      String   @unique
  fileTree  Json?    // Stores FileTree as JSON blob
  leaderId  Int
  createdAt DateTime @default(now())
  users     User[]   @relation("ProjectUsers")
}
\end{lstlisting}

\subsection{API Endpoint Catalogue}

\begin{table}[H]
  \centering
  \caption{Complete REST API endpoint catalogue with authentication status.}
  \label{tab:api}
  \begin{tabular}{lllp{4.5cm}}
    \toprule
    \textbf{Method} & \textbf{Path} & \textbf{Auth} & \textbf{Description} \\
    \midrule
    POST & \texttt{/api/v1/user/register}       & No  & Register new user \\
    POST & \texttt{/api/v1/user/login}           & No  & Login, set cookie \\
    GET  & \texttt{/api/v1/user/logout}          & Yes & Blacklist token \\
    GET  & \texttt{/api/v1/user/profile}         & Yes & Get own profile \\
    GET  & \texttt{/api/v1/user/all}             & Yes & List all other users \\
    POST & \texttt{/api/v1/project/create}       & Yes & Create project \\
    GET  & \texttt{/api/v1/project/getAll}       & Yes & List own projects \\
    GET  & \texttt{/api/v1/project/get-project/:id} & Yes & Get project by ID \\
    POST & \texttt{/api/v1/project/add-user}     & Yes & Add collaborators (leader) \\
    DELETE & \texttt{/api/v1/project/remove-user} & Yes & Remove users (leader) \\
    DELETE & \texttt{/api/v1/project/delete/:id} & Yes & Delete project (leader) \\
    PUT  & \texttt{/api/v1/project/update-file-tree} & Yes & Save file tree \\
    GET  & \texttt{/api/v1/ai/get-result}        & No  & Direct AI query \\
    GET  & \texttt{/health}                      & No  & Health check \\
    \bottomrule
  \end{tabular}
\end{table}
